% Options for packages loaded elsewhere
\PassOptionsToPackage{unicode}{hyperref}
\PassOptionsToPackage{hyphens}{url}
%
\documentclass[
  ignorenonframetext,
]{beamer}
\usepackage{pgfpages}
\setbeamertemplate{caption}[numbered]
\setbeamertemplate{caption label separator}{: }
\setbeamercolor{caption name}{fg=normal text.fg}
\beamertemplatenavigationsymbolsempty
% Prevent slide breaks in the middle of a paragraph
\widowpenalties 1 10000
\raggedbottom
\setbeamertemplate{part page}{
  \centering
  \begin{beamercolorbox}[sep=16pt,center]{part title}
    \usebeamerfont{part title}\insertpart\par
  \end{beamercolorbox}
}
\setbeamertemplate{section page}{
  \centering
  \begin{beamercolorbox}[sep=12pt,center]{part title}
    \usebeamerfont{section title}\insertsection\par
  \end{beamercolorbox}
}
\setbeamertemplate{subsection page}{
  \centering
  \begin{beamercolorbox}[sep=8pt,center]{part title}
    \usebeamerfont{subsection title}\insertsubsection\par
  \end{beamercolorbox}
}
\AtBeginPart{
  \frame{\partpage}
}
\AtBeginSection{
  \ifbibliography
  \else
    \frame{\sectionpage}
  \fi
}
\AtBeginSubsection{
  \frame{\subsectionpage}
}
\usepackage{amsmath,amssymb}
\usepackage{iftex}
\ifPDFTeX
  \usepackage[T1]{fontenc}
  \usepackage[utf8]{inputenc}
  \usepackage{textcomp} % provide euro and other symbols
\else % if luatex or xetex
  \usepackage{unicode-math} % this also loads fontspec
  \defaultfontfeatures{Scale=MatchLowercase}
  \defaultfontfeatures[\rmfamily]{Ligatures=TeX,Scale=1}
\fi
\usepackage{lmodern}
\ifPDFTeX\else
  % xetex/luatex font selection
\fi
% Use upquote if available, for straight quotes in verbatim environments
\IfFileExists{upquote.sty}{\usepackage{upquote}}{}
\IfFileExists{microtype.sty}{% use microtype if available
  \usepackage[]{microtype}
  \UseMicrotypeSet[protrusion]{basicmath} % disable protrusion for tt fonts
}{}
\makeatletter
\@ifundefined{KOMAClassName}{% if non-KOMA class
  \IfFileExists{parskip.sty}{%
    \usepackage{parskip}
  }{% else
    \setlength{\parindent}{0pt}
    \setlength{\parskip}{6pt plus 2pt minus 1pt}}
}{% if KOMA class
  \KOMAoptions{parskip=half}}
\makeatother
\usepackage{xcolor}
\newif\ifbibliography
\setlength{\emergencystretch}{3em} % prevent overfull lines
\providecommand{\tightlist}{%
  \setlength{\itemsep}{0pt}\setlength{\parskip}{0pt}}
\setcounter{secnumdepth}{-\maxdimen} % remove section numbering
% Slide layout defaults for warning-clean scientific decks.
\usepackage{etoolbox}
\IfFileExists{xurl.sty}{\usepackage{xurl}}{}
\IfFileExists{seqsplit.sty}{\usepackage{seqsplit}}{\newcommand{\seqsplit}[1]{#1}}
\protected\def\breakseq#1{\seqsplit{#1}}
\protected\def\breaktt#1{\begingroup\ttfamily\seqsplit{#1}\endgroup}
\setlength{\emergencystretch}{6em}
\tolerance=5000
\hbadness=10000
\hfuzz=1pt
\setlength{\tabcolsep}{2pt}
\AtBeginEnvironment{longtable}{\tiny\renewcommand{\arraystretch}{0.86}\setlength{\tabcolsep}{1pt}}
\AtBeginEnvironment{tabular}{\tiny\renewcommand{\arraystretch}{0.86}\setlength{\tabcolsep}{1pt}}

% Natbib and cross-reference fallbacks — slides don't load natbib
% or cleveref, but combined-PDF manuscript prose may emit these
% commands. The fallback renders citations as a bracketed key list
% and cross-references as detokenized labels so slides don't fail on
% undefined control sequences or raw underscores. \providecommand is
% a no-op if packages load later.
\providecommand{\citep}[1]{[#1]}
\providecommand{\citet}[1]{#1}
\providecommand{\citealp}[1]{#1}
\providecommand{\citeauthor}[1]{#1}
\providecommand{\citeyear}[1]{#1}
\providecommand{\cref}[1]{\texttt{\detokenize{#1}}}
\providecommand{\Cref}[1]{\texttt{\detokenize{#1}}}

\newtheorem{proposition}{Proposition}
\newtheorem{hypothesis}{Hypothesis}
\ifLuaTeX
  \usepackage{selnolig}  % disable illegal ligatures
\fi
\usepackage{bookmark}
\IfFileExists{xurl.sty}{\usepackage{xurl}}{} % add URL line breaks if available
\urlstyle{same}
\hypersetup{
  hidelinks,
  pdfcreator={LaTeX via pandoc}}

\author{}
\date{}

\begin{document}

\section{Multi-Phase Search Strategy and
Results}\label{multi-phase-search-strategy-and-results}

\begin{frame}[fragile,allowframebreaks]{Three-Phase Architecture}
\phantomsection\label{three-phase-architecture}
This review employs a \textbf{three-phase search strategy} that
progressively refines coverage from broad foundational literature
through technology-specific studies to targeted molecular detection
analyses. Each phase builds on prior phases through cross-phase citation
validation and shared deduplication.

\begin{block}{Phase 1: Exoplanet Atmosphere Foundation}
\phantomsection\label{phase-1-exoplanet-atmosphere-foundation}
\textbf{Objective:} Establish the foundational literature on exoplanet
atmospheric studies, covering the broadest range of research from 2010
onwards.

\textbf{Queries:} 1.
\breaktt{"exoplanet\ atmosphere"\ OR\ "exoplanet\ atmospheric"} 2.
\breaktt{"transit\ spectroscopy"\ AND\ "atmosphere"} 3.
\breaktt{"atmospheric\ composition"\ AND\ "exoplanet"}

\textbf{Engines:} arXiv, OpenAlex, Crossref, Semantic Scholar

\textbf{Deterministic Filters:} - Minimum year: 2010 - Maximum year:
2026 - Minimum citation count: 0 (inclusive to capture recent work)

\textbf{Results:} \{\{PHASE\_1\_PAPERS\}\} papers discovered, forming
the foundational knowledge base for the review. This phase captures the
broadest scope of atmospheric research, from observational
characterizations to theoretical modeling studies.
\end{block}

\begin{block}{Phase 2: James Webb Space Telescope Studies}
\phantomsection\label{phase-2-james-webb-space-telescope-studies}
\textbf{Objective:} Identify papers specifically utilizing JWST
observational data for exoplanet atmospheric analysis, capturing the
post-launch (2021+) surge in high-precision spectroscopic observations.

\textbf{Queries:} 1.
\breaktt{"James\ Webb"\ AND\ "atmosphere"\ AND\ ("exoplanet"\ OR\ "planet")}
2. \breaktt{"JWST"\ AND\ "transit"\ AND\ "spectroscopy"} 3.
\breaktt{"NIRSpec"\ OR\ "MIRI"\ OR\ "NIRCam"\ AND\ "exoplanet"}

\textbf{Engines:} arXiv, OpenAlex, Crossref, Semantic Scholar

\textbf{Deterministic Filters:} - Minimum year: 2020 (capturing
pre-launch calibration papers) - Maximum year: 2026

\textbf{Dependencies:} Builds on Phase 1 foundational papers.

\textbf{Results:} \{\{PHASE\_2\_PAPERS\}\} papers after deterministic
filtering. This phase targets records associated with JWST-era
observations and analysis. Any claim about instrumental capability or
scientific findings requires live, source-backed review.
\end{block}

\begin{block}{Phase 3: Molecular Detection Studies}
\phantomsection\label{phase-3-molecular-detection-studies}
\textbf{Objective:} Identify papers focused on detecting and analyzing
specific atmospheric molecules (H₂O, CO₂, CH₄, H₂S, Na, K) in the
configured domain.

\textbf{Queries:} 1.
\breaktt{"water\ vapor"\ AND\ "exoplanet"\ AND\ ("detection"\ OR\ "abundance")}
2.
\breaktt{"carbon\ dioxide"\ AND\ "exoplanet"\ AND\ ("CO2"\ OR\ "atmosphere")}
3. \breaktt{"methane"\ AND\ "exoplanet"\ AND\ ("CH4"\ OR\ "atmosphere")}
4. \breaktt{"hydrogen\ sulfide"\ OR\ "H2S"\ AND\ "exoplanet"} 5.
\breaktt{"sodium"\ OR\ "potassium"\ AND\ "exoplanet\ atmosphere"}

\textbf{Engines:} arXiv, OpenAlex, Crossref, Semantic Scholar

\textbf{Deterministic Filters:} - Minimum year: 2015 - Maximum year:
2026

\textbf{Dependencies:} Builds on both Phase 1 and Phase 2 papers.

\textbf{Results:} \{\{PHASE\_3\_PAPERS\}\} papers covering molecular
detection across diverse exoplanet populations, from hot Jupiters to
terrestrial worlds.
\end{block}
\end{frame}

\begin{frame}[allowframebreaks]{Cross-Phase Analysis}
\phantomsection\label{cross-phase-analysis}
\begin{block}{Phase Overlap}
\phantomsection\label{phase-overlap}
The three phases produced complementary but overlapping corpora, with
\{\{CROSS\_PHASE\_OVERLAP\_PCT\}\}\% average Jaccard similarity between
phases. Papers discovered in multiple phases are tracked with full
provenance, enabling analysis of papers that span multiple research
paradigms.
\end{block}

\begin{block}{Citation Validation}
\phantomsection\label{citation-validation}
Cross-phase citation analysis records that
\{\{CROSS\_PHASE\_CITATION\_RATE\}\}\% of later-phase papers cite
foundational work from Phase 1. This is a descriptive connection
statistic for the retained corpus; it does not establish methodological
coherence across the field.
\end{block}

\begin{block}{Combined Corpus}
\phantomsection\label{combined-corpus}
The deduplicated combined corpus contains \{\{CORPUS\_SIZE\}\} unique
papers spanning \{\{YEAR\_START\}\}--\{\{YEAR\_END\}\}
(\{\{YEAR\_SPAN\}\} years). It provides a multi-phase description of the
retained retrieval slice, not a field-wide account.
\end{block}
\end{frame}

\begin{frame}[fragile,allowframebreaks]{LLM-Based Content Filtering}
\phantomsection\label{llm-based-content-filtering}
Three LLM-based content filters were designed for this review:

\begin{enumerate}
\item
  \textbf{Study Type Classification}: Categorizes papers as
  observational, theoretical, review, or other. Keeps only observational
  and theoretical studies for analysis.
\item
  \textbf{JWST Data Analysis}: Identifies papers that analyze actual
  JWST observational data, distinguishing from theoretical JWST
  predictions.
\item
  \textbf{Molecular Detection Focus}: Filters for papers primarily
  focused on detecting and measuring specific atmospheric molecules.
\end{enumerate}

These filters are configurable through \breaktt{manuscript/config.yaml}
and can be enabled for specific phases or applied across the entire
corpus. When enabled, they use a local Ollama LLM instance for
cost-effective, privacy-preserving content classification.
\end{frame}

\end{document}
